\chapter{局限性探讨与总结}

综合上述数学推导，本文在数据非独立同分布且无中心服务器的网络环境中，构建了一套抵御拜占庭故障的多智能体优化理论框架。通过引入 $(\tilde{\beta}, \tilde{\gamma})$-可容许性度量，本文在底层机制上重构了协同优化的物理边界，并在完全图拓扑条件下，实现了与理论容错能力界限（$\tilde{\gamma} \le n-\phi-b$）相匹配的系统收敛性。

该研究目前仍面临若干有待解决的问题。首先，现有的收敛性分析（如式 \ref{eq:dist_bound} 所示）仅确立了时间趋于无穷大时的渐近收敛性，未能在残差项积分中显式推导出有限时间收敛率。缺乏收敛率界限限制了该算法在对实时性要求较高的计算场景中的性能评估。其次，在分布式算法的具体实现上，为构造严格的马尔可夫演化方程并满足引理 \ref{lem:3} 的代数等价性，当前算法要求节点在每次通信迭代中同时交换状态估计与局部梯度序列。这一双重信息交换机制增加了系统整体的通信带宽开销。未来探寻如何在仅交换状态变量的前提下，利用一阶梯度的历史差分完成等效的几何边界约束，是推进该领域实际应用的重要研究方向。
